\chapter{Symmetries}

\ex{
    In physics symmetries are often first introduced in the context of classical mechanics in order to simplify computations.
}
\defn{Symmetry}{
    A \textbf{symmetry} is an invertible map
    \[
        s : \mathbb{R}^n \to \mathbb{R}^n
    \]
    such that, whenever \(x(t)\) is a solution of the equations of motion, also
    \[
        s \circ x(t)
    \]
    is a solution.
}

\ex{
    Consider a free particle, so that \(\vec{F} = 0\). Then the equations of motion
    have solutions
    \[
        x(t) = x_0 + t v_0,
        \qquad x_0, v_0 \in \mathbb{R}^n.
    \]
    For any \(T \in \mathbb{R}^n\), the translation
    \[
        s_T : \mathbb{R}^n \to \mathbb{R}^n,
        \qquad x \mapsto x + T
    \]
    is a symmetry.
}

\ex{
    In quantum mechanics, a symmetry is typically a unitary operator
    \[
        S : \mathcal{H} \to \mathcal{H}
    \]
    on the Hilbert space \(\mathcal{H}\) such that
    \[
        [S,H] = 0,
    \]
    where \(H : \mathcal{H} \to \mathcal{H}\) is the Hamiltonian.

    For
    \[
        \mathcal{H} = L^2(\mathbb{R}),
        \qquad
        H = -\frac{1}{2m}\partial_x^2,
    \]
    the translation operator
    \[
        S_T : L^2(\mathbb{R}) \to L^2(\mathbb{R}),
        \qquad
        (S_T f)(x) = f(x+T),
    \]
    and the reflection operator
    \[
        S_R : L^2(\mathbb{R}) \to L^2(\mathbb{R}),
        \qquad
        (S_R f)(x) = f(-x),
    \]
    are symmetries.
}

\begin{center}
\begin{tabular}{p{0.42\textwidth}|p{0.42\textwidth}}
\textbf{Abstract structure} & \textbf{Concrete realization} \\[0.5em]
\hline

\( \text{group } G \)
&
\( \text{representation } \rho_V \)
\\[1em]

\( g \in G \)
&
\( \rho(g) : V \to V \)
\\[1em]

\( e \in G \quad (\text{``do nothing''}) \)
&
\( \rho(e) = \mathrm{id}_V \)
\\[1em]

\( h \circ g \in G \quad (\text{symmetry composition}) \)
&
\( \rho(h \circ g) = \rho(h)\,\rho(g) \)
\\[1em]

\(
\exists\, g^{-1} \text{ s.t. } g^{-1}\circ g = g \circ g^{-1} = e
\)
&
\(
\rho(g^{-1})\,\rho(g)
=
\rho(g)\,\rho(g^{-1})
=
\mathrm{id}_V
\)
\end{tabular}
\end{center}

\defn{Group}{
    A \textbf{group} is a set \(G\) together with a map
    \[
        \circ : G \times G \to G,
        \qquad
        (g,h) \mapsto g \circ h,
    \]
    such that:
    \begin{enumerate}
        \item \(\circ\) is associative, that is,
        \[
            (g \circ h) \circ k = g \circ (h \circ k)
            \qquad \forall g,h,k \in G.
        \]
        \item There exists an element \(e \in G\) such that
        \[
            e \circ g = g \circ e = g
            \qquad \forall g \in G.
        \]
        \item For every \(g \in G\), there exists an element \(g^{-1} \in G\) such that
        \[
            g \circ g^{-1} = g^{-1} \circ g = e.
        \]
    \end{enumerate}
}

\corp{
    In a group, the unit element and inverses are unique.
}{
    Suppose \(e\) and \(e'\) are both unit elements. Then
    \[
        e = e \circ e' = e'.
    \]
    Hence the unit is unique.

    Now let \(g \in G\), and suppose \(h\) and \(h'\) are both inverses of \(g\). Then
    \[
        h = h \circ e = h \circ (g \circ h') = (h \circ g) \circ h' = e \circ h' = h'.
    \]
    Hence the inverse of any element is unique.
}{}

\defn{Representation}{
    Let \((G,\circ)\) be a group. Usually we simply write \(G\) and leave the map
    \(\circ\) implicit. A \textbf{representation} of \(G\) is a vector space \(V\)
    together with a map
    \[
        \rho : G \to \operatorname{End}(V)
    \]
    such that:
    \begin{enumerate}
        \item
        \( \rho(e) = \mathrm{id}_V. \)
        \item
        \( \rho(g)\rho(h) = \rho(g \circ h)
            \qquad \forall g,h \in G. \)
    \end{enumerate}
}

We will study these two definitions in this course. The basic structure of the course is as follows:

\begin{center}
\resizebox{0.92\textwidth}{!}{
\begin{tikzpicture}[
    x=1cm, y=1cm,
    >=Stealth,
    every node/.style={align=center},
    box/.style={inner sep=2pt},
    arr/.style={thick, ->, rounded corners=10pt}
]

\node[box] (groups) at (0,6.2) {\underline{\Large Groups}};

\node[box] (disc) at (-4.7,4.4) {\underline{\Large Discrete groups}};
\node[box] (cont) at (4.7,4.4) {\underline{\Large Continuous Groups}};

\coordinate (groupsplit) at (0,5.45);
\draw[arr] (groups.south) -- (groupsplit) -| (disc.north);
\draw[arr] (groups.south) -- (groupsplit) -| (cont.north);

\node[anchor=west] at (-6.7,3.65) {\large - e.g.\ Reflections of polygons};

\draw[thick] (-5.9,1.2) -- (-4.9,2.932) -- (-3.9,1.2) -- cycle;
\draw[thick] (-4.9,0.5) -- (-4.9,3.3);
\draw[thick] (-6.2,1.0268) -- (-3.6,2.5292);
\draw[thick] (-6.2,2.5292) -- (-3.6,1.0268);
\node at (-3.3,1.7) {\Large $S_3$};

\node[box] (finite) at (-6.0,-1.1) {\underline{\Large finite groups}};
\node[box] (infinite) at (-2.0,-1.1) {\underline{\Large infinite groups}};

\coordinate (discsplit) at (-4.9,0.3);
\draw[arr] (discsplit) -| (finite.north);
\draw[arr] (discsplit) -| (infinite.north);

\node[anchor=west] at (-7.3,-2.2) {\Large $|G| < \infty$};
\node[anchor=west] at (-7.3,-3.0) {\large - e.g.\ $S_3$};

\node at (-2.0,-2.3) {\Large $(\mathbb{Z},+)$};

\node[anchor=west] at (2.8,3.3) {\Large Lie groups};
\node[anchor=west] at (2.8,2.4) {\large - e.g.\ Translations};
\node[anchor=west] at (2.8,1.5) {\large - e.g.\ $O(1,3)$ Lorentz transf.};

\node at (5.0,0.4) {\Large $\Downarrow$};

\node[box] at (4.9,-0.5) {\underline{\Large Lie algebra}};
\node[anchor=west] at (3.2,-1.6) {\large - infinitesimal sym.};

\end{tikzpicture}
}
\end{center}
